\chapter{Supplementary Materials}
\label{app:supplementary}

This appendix provides detailed specifications, configurations, and reference materials supporting the experimental work presented in Chapters~\ref{ch:rag}--\ref{ch:strategies}. Section~\ref{app:model_specs} catalogs all evaluated models; Section~\ref{app:hardware} details hardware platforms; Section~\ref{app:exp_configs} presents experimental configurations; Section~\ref{app:invoice_errors} illustrates multi-agent error patterns; Section~\ref{app:prompts} documents prompt templates; and Section~\ref{app:metrics} consolidates the seven novel evaluation metrics introduced throughout this dissertation.

\section{Complete Model Specifications}
\label{app:model_specs}

Table~\ref{tab:full_model_specs} provides detailed specifications for all models evaluated in this dissertation, organized by model family. Models are grouped into proprietary alternatives (used as comparison baselines), reasoning-specialized architectures (central to Chapter~\ref{ch:strategies}), and general-purpose open-source models (evaluated across all three research pillars).

\begin{longtable}{@{}p{3.5cm}p{1.8cm}p{2cm}p{2.2cm}p{3cm}@{}}
\caption{Complete Model Specifications} \label{tab:full_model_specs} \\
\toprule
\textbf{Model} & \textbf{Parameters} & \textbf{Context} & \textbf{Architecture} & \textbf{Notes} \\
\midrule
\endfirsthead
\multicolumn{5}{c}{{\tablename\ \thetable{} -- continued from previous page}} \\
\toprule
\textbf{Model} & \textbf{Parameters} & \textbf{Context} & \textbf{Architecture} & \textbf{Notes} \\
\midrule
\endhead
\midrule
\multicolumn{5}{r}{{Continued on next page}} \\
\endfoot
\bottomrule
\endlastfoot
\multicolumn{5}{c}{\textit{Proprietary Models (Comparison Baselines)}} \\
\midrule
GPT-4 & Undisclosed & 128K & MoE (speculated) & OpenAI; RAG baseline \\
GPT-4o & Undisclosed & 128K & Multimodal & OpenAI; sentiment/reasoning \\
GPT-4o-mini & Undisclosed & 128K & Multimodal & OpenAI; cost-optimized \\
GPT-4.1 & Undisclosed & 1M & Multimodal & OpenAI, Apr 2025; agentic \\
GPT-4.1-nano & Undisclosed & 1M & Multimodal & OpenAI; edge annotation \\
GPT-4.1-mini & Undisclosed & 1M & Multimodal & OpenAI; LAJ framework \\
Claude Sonnet 4 & Undisclosed & 200K & Decoder-only & Anthropic; agentic eval \\
\midrule
\multicolumn{5}{c}{\textit{DeepSeek Family}} \\
\midrule
DeepSeek-V3 & 671B & 128K & MoE & Base model for DSR1 \\
DeepSeek-R1 (Full) & 671B & 128K & MoE & RL-trained reasoning \\
DeepSeek-R1:8B & 8B & 64K & Decoder-only & Distilled from Llama3.1-8B \\
DeepSeek-R1:14B & 14B & 64K & Decoder-only & Distilled from Qwen2.5-14B \\
DeepSeek-R1:32B & 32B & 64K & Decoder-only & Distilled from Qwen2.5-32B \\
DeepSeek-R1:70B & 70B & 64K & Decoder-only & Distilled from Llama3.3-70B \\
\midrule
\multicolumn{5}{c}{\textit{Qwen Family}} \\
\midrule
Qwen2-7B & 7B & 128K & Decoder-only & Translation experiments \\
Qwen2-72B & 72B & 128K & Decoder-only & Translation; fine-tuning \\
Qwen2.5-0.5B & 0.5B & 128K & Decoder-only & Edge deployment \\
Qwen2.5-1.5B & 1.5B & 128K & Decoder-only & Edge deployment \\
Qwen2.5-3B & 3B & 128K & Decoder-only & Edge deployment \\
Qwen2.5-7B & 7B & 128K & Decoder-only & 18T tokens; RAG/edge \\
Qwen2.5-14B & 14B & 128K & Decoder-only & Base for DSR1-14B \\
Qwen2.5-32B & 32B & 128K & Decoder-only & Base for DSR1-32B; agentic \\
Qwen2.5-72B & 72B & 128K & Decoder-only & 18T tokens; reasoning \\
Qwen2.5-Plus & API & 128K & Decoder-only & Soccer betting sentiment \\
Qwen3-0.6B & 0.6B & 32K & Decoder-only & T/N modes \\
Qwen3-1.7B & 1.7B & 32K & Decoder-only & T/N modes \\
Qwen3-4B & 4B & 32K & Decoder-only & T/N modes \\
Qwen3-8B & 8B & 32K & Decoder-only & T/N modes; best Qwen3 \\
Qwen3-14B & 14B & 32K & Decoder-only & T/N modes \\
Qwen3-30B & 30B & 32K & Decoder-only & T/N modes \\
Qwen3-32B & 32B & 32K & Decoder-only & T/N modes; annotation \\
\midrule
\multicolumn{5}{c}{\textit{Llama Family}} \\
\midrule
Llama-2-7B & 7B & 4K & Decoder-only & 2T tokens \\
Llama-2-13B & 13B & 4K & Decoder-only & 2T tokens \\
Llama-2-70B & 70B & 4K & Decoder-only & 2T tokens \\
Llama-3-8B & 8B & 8K & Decoder-only & 15T tokens \\
Llama-3-70B & 70B & 8K & Decoder-only & 15T tokens \\
Llama-3.1-8B & 8B & 128K & Decoder-only & Base for DSR1-8B; edge \\
Llama-3.1-70B & 70B & 128K & Decoder-only & 15T+ tokens; fine-tuning \\
Llama-3.2-1B & 1B & 128K & Decoder-only & Edge-optimized \\
Llama-3.2-3B & 3B & 128K & Decoder-only & Edge-optimized \\
Llama-3.2-vision:11B & 11B & 128K & Multimodal & Vision-language; edge \\
Llama-3.2-vision:90B & 90B & 128K & Multimodal & Vision-language \\
Llama-3.3-70B & 70B & 128K & Decoder-only & Base for DSR1-70B \\
\midrule
\multicolumn{5}{c}{\textit{Reasoning-Specialized Models}} \\
\midrule
Granite3.3-2B & 2B & 128K & Decoder-only & IBM; T/N modes \\
Granite3.3-8B & 8B & 128K & Decoder-only & IBM; T/N modes \\
Magistral-24B & 24B & 128K & Decoder-only & Mistral; RL reasoning \\
GPT-OSS:20B (low) & 20B & 32K & Decoder-only & Minimal reasoning (RI$<$0.5) \\
GPT-OSS:20B (medium) & 20B & 32K & Decoder-only & Balanced reasoning \\
GPT-OSS:20B (high) & 20B & 32K & Decoder-only & Intensive reasoning (RI$>$0.85) \\
GPT-OSS:120B (low) & 120B & 32K & Decoder-only & Minimal reasoning \\
GPT-OSS:120B (medium) & 120B & 32K & Decoder-only & Balanced; best F1 (86.2\%) \\
GPT-OSS:120B (high) & 120B & 32K & Decoder-only & Intensive reasoning \\
\midrule
\multicolumn{5}{c}{\textit{Other Open-Source Models}} \\
\midrule
Orca-2-7B & 7B & 4K & Decoder-only & Synthetic training \\
Orca-2-13B & 13B & 4K & Decoder-only & Synthetic; RAG champion \\
Mistral-7B & 7B & 32K & Decoder-only & Sliding window attn. \\
Phi-3.5-mini & 3.8B & 128K & Decoder-only & Synthetic training \\
Gemma-2B & 2B & 8K & Decoder-only & Google; 2T tokens \\
Gemma-7B & 7B & 8K & Decoder-only & Google; 6T tokens \\
InternLM2.5-7B & 7B & 32K & Decoder-only & Shanghai AI Lab \\
InternLM2.5-7B-1M & 7B & 1M & Decoder-only & Extended context; reasoning \\
InternLM2.5-20B & 20B & 32K & Decoder-only & Shanghai AI Lab \\
Functionary-small & 8B & 32K & Function-calling & Tool-optimized; agentic \\
Functionary-medium & 70B & 32K & Function-calling & Tool-optimized; agentic \\
\end{longtable}

\textbf{Notes:} 
\begin{itemize}[leftmargin=*]
    \item T/N modes = Thinking/Non-thinking modes with runtime-activated reasoning (Section~\ref{subsec:task_complexity})
    \item MoE = Mixture of Experts architecture
    \item RI = Reasoning Intensity metric introduced in Section~\ref{subsec:reasoning_deployment}
    \item Context lengths represent maximum supported; actual usage may vary by deployment configuration
    \item Models marked ``edge'' were evaluated across consumer hardware platforms (Section~\ref{subsec:ollama_edge})
    \item Models marked ``agentic'' were evaluated in multi-agent workflows (Chapter~\ref{ch:agentic})
\end{itemize}

\section{Hardware Specifications}
\label{app:hardware}

Table~\ref{tab:hardware_specs} details the six primary hardware platforms used for evaluation in this dissertation, spanning enterprise GPUs, consumer hardware, and edge devices. These platforms correspond to the edge deployment evaluation in Section~\ref{subsec:ollama_edge}.

\begin{table}[htbp]
\centering
\caption{Hardware Platforms for Model Evaluation}
\label{tab:hardware_specs}
\resizebox{\textwidth}{!}{%
\begin{tabular}{@{}p{3.5cm}p{9cm}@{}}
\toprule
\textbf{Platform} & \textbf{Specifications} \\
\midrule
\multicolumn{2}{c}{\textit{Enterprise GPU}} \\
\midrule
NVIDIA H100 & 96GB HBM3 VRAM, 16,896 CUDA cores, 528 Tensor cores, PCIe 5.0, 3.9 TB/s bandwidth, Ubuntu 24.04.1 LTS, CUDA 12.4, Ollama 0.11.4 \\
\midrule
\multicolumn{2}{c}{\textit{Professional Workstation}} \\
\midrule
NVIDIA RTX A6000 Ada & 48GB GDDR6 ECC VRAM, 18,176 CUDA cores, 568 Tensor cores, AMD Threadripper PRO 5975WX (32 cores), 256GB DDR4-3200 ECC, Ubuntu 22.04 LTS / Windows 11 Pro + WSL2 \\
\midrule
\multicolumn{2}{c}{\textit{Consumer Desktop}} \\
\midrule
NVIDIA RTX 4090 & 24GB GDDR6X VRAM, 16,384 CUDA cores, 512 Tensor cores, Intel Core i9-13900K (24 cores), 64GB DDR5-5600, Ubuntu 22.04 LTS via WSL2, CUDA 12.1, Ollama 0.5.4--0.6.8 \\
\midrule
\multicolumn{2}{c}{\textit{Consumer Laptop}} \\
\midrule
NVIDIA RTX 4090 Laptop & 16GB GDDR6 VRAM, 9,728 CUDA cores, 304 Tensor cores, Intel Core i9-14900HX, 64GB DDR5, Windows 11 Pro / WSL2 Ubuntu 22.04 (\textbf{21$\times$ speedup over native Windows for small models}) \\
\midrule
\multicolumn{2}{c}{\textit{Apple Silicon}} \\
\midrule
M3 Max MacBook Pro & 96GB Unified Memory, 16-core CPU, 40-core GPU, 16-core Neural Engine, macOS Sequoia 15.4, MLX 0.12.0, Ollama 0.5.4 \\
M4 Max MacBook Pro & 48GB Unified Memory, 16-core CPU, 40-core GPU, 16-core Neural Engine, macOS Sequoia 15.1, Ollama 0.5.4 \\
\midrule
\multicolumn{2}{c}{\textit{Edge Device}} \\
\midrule
NVIDIA Jetson AGX Orin & 64GB LPDDR5 Memory, 2,048 CUDA cores, 64 Tensor cores, 12-core Arm Cortex-A78AE CPU, up to 275 TOPS INT8, JetPack 6.0, Ubuntu 20.04 \\
\bottomrule
\end{tabular}
}
\end{table}

\subsection{Platform Selection Rationale}

The hardware platforms were selected to represent the full spectrum of enterprise deployment scenarios addressed in Chapter~\ref{ch:strategies}:

\begin{itemize}[leftmargin=*]
    \item \textbf{H100:} Cloud/data center deployments requiring maximum throughput; primary platform for 504-configuration task complexity evaluation (Section~\ref{subsec:task_complexity})
    \item \textbf{RTX A6000 Ada:} Professional workstations for development and medium-scale inference; energy efficiency evaluation (Figure~\ref{fig:energy_efficiency})
    \item \textbf{RTX 4090 Desktop:} Cost-effective on-premises deployment with strong performance; multi-agent invoice reconciliation (Chapter~\ref{ch:agentic}), achieving 99.90\% success rate with qwen2.5:32b at 523 Joules per task
    \item \textbf{RTX 4090 Laptop:} Mobile professional use cases; 16GB VRAM constraint defines viable reasoning models (qwen3:8b optimal at 84.8\% F1, 13GB VRAM)
    \item \textbf{Apple Silicon:} macOS-native deployment and unified memory advantages; competitive edge performance
    \item \textbf{Jetson AGX Orin:} Edge deployment for privacy-sensitive and latency-critical applications
\end{itemize}

\subsection{Key Hardware Findings Summary}

The following hardware-specific findings from Section~\ref{subsec:ollama_edge} inform deployment decisions:

\begin{enumerate}[leftmargin=*]
    \item \textbf{WSL Performance:} Windows Subsystem for Linux delivers up to 21$\times$ throughput improvement over native Windows for smaller models (qwen2.5:0.5b: 16,933 vs. 816 tokens/s on RTX A6000)
    \item \textbf{VRAM Constraints:} Consumer 16GB GPUs support models up to 14B parameters for reasoning tasks; 24GB enables 32B deployment
    \item \textbf{Cross-Platform Stability:} Larger models exhibit increased consistency (F1 std: 0.07--0.23\%) compared to small models (0.26--0.90\%)
    \item \textbf{Energy Efficiency:} Small models (1.5B--3B) achieve highest performance-to-power ratios (0.161--0.163 \%/J); large models consume 1,283--1,876 Joules per inference
\end{enumerate}

\section{Detailed Experimental Configurations}
\label{app:exp_configs}

\subsection{RAG Experiments Configuration (Chapter~\ref{ch:rag})}

\begin{verbatim}
# Embedding Model Configuration
embedding_model: hkunlp/instructor-large
embedding_dim: 768
max_seq_length: 512

# Retrieval Configuration
vector_store: FAISS
index_type: IVF_FLAT
num_clusters: 100
top_k_retrieval: 10

# Chunking Configuration
chunk_size: 500  # tokens
chunk_overlap: 50  # tokens
separator: "\n\n"

# Generation Configuration
temperature: 0.7
max_new_tokens: 1024
top_p: 0.95
repetition_penalty: 1.0-1.2  # varied in experiments

# RAP Metric Configuration
rap_exponent: 3  # cubic penalty validated via ablation
performance_metric: RAGAS_composite  # or task-specific F1

# Hardware
gpu: NVIDIA RTX 4090 (24GB) / A6000 (48GB)
precision: float16
\end{verbatim}

\subsection{Fine-Tuning Configuration (Chapters~\ref{ch:rag}, \ref{ch:strategies})}

\begin{verbatim}
# LoRA Configuration
lora_r: 64
lora_alpha: 128
lora_dropout: 0.1
target_modules: ["q_proj", "v_proj", "k_proj", "o_proj"]

# Training Configuration
learning_rate: 2e-4
batch_size: 8
gradient_accumulation_steps: 4
num_epochs: 3  # with early stopping at 0.2 epoch granularity
warmup_ratio: 0.1
weight_decay: 0.01
optimizer: AdamW

# QLoRA Specific (for models >= 20B)
quantization_bits: 4
bnb_4bit_compute_dtype: float16
bnb_4bit_quant_type: nf4
bnb_4bit_use_double_quant: True

# Early Stopping (per Section 5.3.1.5)
# Larger models peak earlier: 0.8 epochs (72B) vs 1.0 epochs (7B)
# VCR improves from 1.17% to 100% within 0.2 epochs
\end{verbatim}

\subsection{Multi-Agent System Configuration (Chapter~\ref{ch:agentic})}

\begin{verbatim}
# LangGraph Configuration
graph_type: StateGraph
checkpoint_store: MemorySaver
recursion_limit: 25

# Agent Configuration
max_iterations: 10
tool_timeout: 30  # seconds
retry_attempts: 3

# Database Configuration
db_type: SQLite
connection_mode: in-memory

# Email Agent Configuration
imap_server: imap.gmail.com
smtp_server: smtp.gmail.com
fetch_limit: 50

# Ollama Configuration (Edge Deployment)
ollama_version: 0.6.8
quantization: Q4_K_M
context_length: 8192

# HMASP Four-Level Hierarchy
hierarchy_levels: [orchestrator, domain_supervisors, 
                   specialized_agents, tools]
shared_state_variables: true
decoupled_message_states: true
structured_handoff_protocols: true

# Success Metrics
target_ASR: 0.93  # minimum for production-ready
target_TSR: 0.97  # task success rate
ECR_retry_budget: 0.15  # 15% overhead for reliability failures
\end{verbatim}

\subsection{Task Complexity Experiments Configuration (Section~\ref{subsec:task_complexity})}

The following configuration was used for the 504-configuration evaluation of reasoning effectiveness across sentiment analysis tasks.

\begin{verbatim}
# Model Families Evaluated (7 families)
deepseek_r1: [671B-full, 8B, 14B, 32B, 70B]
deepseek_v3: [671B]  # Base model comparison
llama_bases: [3.1-8B, 3.3-70B]
qwen25_bases: [14B, 32B]
qwen3: [4B, 8B, 14B, 30B, 32B]  # T/N modes
granite33: [2B, 8B]  # T/N modes
magistral: [24B]  # T/N modes

# Datasets (3 complexity levels)
imdb:
  task: binary_sentiment
  classes: [Positive, Negative]
  test_size: 1000
  complexity: simple
  expected_reasoning_benefit: negative  # -4.8 pp avg

amazon_reviews:
  task: 5class_sentiment
  classes: [Strongly Negative, Negative, Neutral, 
            Positive, Strongly Positive]
  test_size: 1838  # human-annotated
  complexity: moderate
  expected_reasoning_benefit: mixed  # -3.6 pp avg

goemotions:
  task: 27class_emotion
  classes: [admiration, amusement, anger, annoyance,
            approval, caring, confusion, curiosity,
            desire, disappointment, disapproval, disgust,
            embarrassment, excitement, fear, gratitude,
            grief, joy, love, nervousness, optimism,
            pride, realization, relief, remorse, 
            sadness, surprise]
  test_size: 1000
  complexity: high
  expected_reasoning_benefit: positive  # +2.0 pp avg
  note: single-label subset only

# Few-Shot Configuration (7 levels)
shot_levels: [0, 5, 10, 20, 30, 40, 50]
sampling_strategy: stratified
random_seed: 42

# Inference Configuration
ollama_version: 0.6.8
quantization: Q4_K_M  # MXFP4 for GPT-OSS
context_length: 8192
temperature: 0.0  # deterministic
max_tokens: 256

# Hardware
primary: NVIDIA H100 (96GB)
validation: NVIDIA RTX 4090 (24GB)
cross_platform: [M3 Max, RTX 4090 Laptop]  # temporal robustness

# Metrics Computed
performance: [F1, Accuracy, Weighted-F1]
efficiency: [mean_latency_per_sample, VRAM_usage, speed_ratio]
reasoning: [MRT, MOT, RI]  # Mean Reasoning/Output Tokens

# Pareto Frontier Analysis
pareto_axes: [F1_score, log_latency]
frontier_identification: non-dominated_sorting
\end{verbatim}

\textbf{Configuration Notes:}
\begin{itemize}[leftmargin=*]
    \item Total configurations: 7 model families $\times$ 3 datasets $\times$ 7 shot levels $\times$ T/N variants = 504
    \item DeepSeek-R1 (full) and DeepSeek-V3 accessed via official API; all others via Ollama
    \item Thinking (T) and Non-thinking (N) modes evaluated separately for Qwen3, Granite3.3, and Magistral
    \item Stratified sampling ensures balanced class representation in few-shot exemplars
    \item Failure rates by dataset: IMDB 100\%, Amazon 80\%, GoEmotions 50\% (reasoning underperforms)
\end{itemize}

\subsection{Reasoning Deployment Experiments Configuration (Section~\ref{subsec:reasoning_deployment})}

\begin{verbatim}
# Model Families (17 variants)
deepseek_r1: [8B, 14B, 32B, 70B]
qwen3: [0.6B, 1.7B, 4B, 8B, 14B, 30B, 32B]
gpt_oss: [20B-low, 20B-medium, 20B-high,
          120B-low, 120B-medium, 120B-high]

# Dataset
amazon_reviews_5class:
  train_split: 70%
  test_split: 30%
  total_samples: 1838

# Temporal Robustness Evaluation
testset2:
  samples: 574
  collection_date: mid-2025
  annotation: manual

# Quantization
deepseek_qwen: Q4_K_M
gpt_oss: MXFP4

# Ollama Configuration
version: 0.11.4
context_length: 8192

# Reasoning Intensity (RI) Regimes
minimal: RI < 0.5      # GPT-OSS low configurations
balanced: 0.5 <= RI <= 0.85  # optimal regime
intensive: RI > 0.85   # DeepSeek-R1 models

# Consumer Hardware Constraint
rtx_4090_laptop_vram: 16GB
viable_models: [qwen3:8b (13GB), qwen3:4b (11GB), 
                qwen3:1.7b (6.8GB)]
optimal_choice: qwen3:8b  # 84.8% F1, 8.35s inference
\end{verbatim}

\subsection{Soccer Betting Experiments Configuration (Section~\ref{subsec:soccer_betting})}

\begin{verbatim}
# Data Sources
match_statistics: API-Football
seasons: 2015-16 to 2023-24 (9 seasons)
betting_odds: Football-Data.co.uk (BetFair 1X2)
commentaries: BBC Sport (Python Selenium scraping)

# Data Partitioning
feature_selection: 2015-19
training: 2015-22
testing: 2022-23
betting_simulation: 2023-24

# LLM Sentiment Analysis
models: [DeepSeek-R1, GPT-4o-mini, GPT-4.1, Qwen2.5-Plus]
prompting: zero-shot
sentiment_scale: 1-5 (Strongly Negative to Strongly Positive)
feature_computation: home_sentiment - away_sentiment

# Machine Learning Models
classifiers: [LogisticRegression, SVC, RandomForest, XGBoost]
tuning_strategies: [accuracy-tuned, F1-tuned]
hyperparameter_search: GridSearchCV
cross_validation: 5-fold

# Betting Simulation
initial_bankroll: $10,000
kelly_fraction: 0.125  # one-eighth Kelly
bet_condition: positive_expected_value

# Key Results Reference
best_return: 279.48%  # Random Forest + GPT-4.1 (F1-tuned)
second_best: 247.97%  # XGBoost Baseline (F1-tuned)
third_best: 174.01%   # Logistic Regression + DeepSeek-R1 (F1-tuned)
\end{verbatim}

\section{Invoice Reconciliation Error Examples}
\label{app:invoice_errors}

The following examples illustrate the three outcome categories tracked by the multi-agent invoice reconciliation system (Chapter~\ref{ch:agentic}). These examples demonstrate the diagnostic value of the structured error taxonomy and ASR metric introduced in Section~\ref{sec:diagnostic_framework}.

\subsection{SUCCESS Example}

\begin{verbatim}
Input: Reconcile invoice INV-2024-001 for customer ABC Corp
       with recent email communications.

Agent Trace:
1. Data Engineering Agent: Retrieved invoice details
   - Invoice ID: INV-2024-001
   - Amount: $15,000
   - Status: Unpaid
   - Customer: ABC Corp
   - ASR Step 1: PASS (tool initialization)
   - ASR Step 2: PASS (parameter handling)
   - ASR Step 3: PASS (query execution)

2. Email Agent: Processed 3 relevant emails
   - Email 1: Payment confirmation mention
   - Email 2: Wire transfer receipt attachment
   - Attachment OCR: Amount $15,000, Ref: INV-2024-001
   - ASR Step 4: PASS (email retrieval)
   - ASR Step 5: PASS (OCR extraction)

3. Reconciliation Agent: Match confirmed
   - Invoice amount matches payment
   - Reference numbers align
   - Status updated to: Paid
   - ASR Step 6: PASS (matching logic)
   - ASR Step 7: PASS (status update)

Output:
{
  "status": "SUCCESS",
  "invoice_id": "INV-2024-001",
  "amount": 15000,
  "payment_matched": true,
  "reconciliation_date": "2024-11-15T10:30:00Z",
  "asr_score": 1.0,  # 7/7 steps successful
  "energy_consumed_joules": 523
}
\end{verbatim}

\subsection{NO INVOICE Error Example}

\begin{verbatim}
Input: Reconcile invoice INV-2024-999 for customer XYZ Inc.

Agent Trace:
1. Data Engineering Agent: Query executed
   - SQL: SELECT * FROM invoices WHERE id = 'INV-2024-999'
   - Result: No rows returned
   - ASR Step 1: PASS (tool initialization)
   - ASR Step 2: PASS (parameter handling)
   - ASR Step 3: PASS (query execution - valid empty result)

2. Email Agent: Searched for related emails
   - Found 0 matching emails
   - ASR Step 4: PASS (search completed)

3. Reconciliation Agent: Unable to proceed
   - No invoice data available
   - No email evidence found
   - ASR Step 5: PASS (correct termination)

Output:
{
  "status": "NO_INVOICE",
  "error": "Invoice INV-2024-999 not found in database",
  "invoice_id": "INV-2024-999",
  "suggestion": "Verify invoice ID or check if invoice exists",
  "asr_score": 1.0,  # All executed steps successful
  "error_category": "Data Availability"  # Not a system error
}
\end{verbatim}

\subsection{ERROR Example (Tool Invocation Failure)}

\begin{verbatim}
Input: Process payment reconciliation for all outstanding invoices

Agent Trace:
1. Data Engineering Agent: Attempted batch query
   - Error: Query timeout after 30 seconds
   - Reason: Too many records requested
   - ASR Step 1: PASS (tool initialization)
   - ASR Step 2: FAIL (parameter handling - unbounded query)
   - Error Category: "Parameter Validation" (from structured error taxonomy)

2. System: Agent retry triggered (attempt 2/3)
   - Error persisted
   - ASR Step 3: FAIL (retry unsuccessful)

3. Reconciliation Agent: Task failed
   - Unable to retrieve invoice data
   - Workflow terminated
   - ASR Step 4: PASS (graceful termination)

Output:
{
  "status": "ERROR",
  "error_type": "QueryTimeout",
  "error_category": "Parameter Validation",
  "message": "Unable to process batch request",
  "suggestion": "Process invoices individually or in smaller batches",
  "asr_score": 0.5,  # 2/4 steps successful
  "retry_count": 2,
  "ecr_impact": "First-attempt failure; retry overhead incurred"
}
\end{verbatim}

\subsection{Error Taxonomy Reference}

The structured error taxonomy introduced in Chapter~\ref{ch:agentic} classifies failures as follows:

\begin{enumerate}[leftmargin=*]
    \item \textbf{Tool Initialization Errors:} Agent fails to properly initialize required tools
    \item \textbf{Parameter Validation Errors:} Invalid or unbounded parameters passed to tools
    \item \textbf{Parameter Type Errors:} Type mismatches in tool arguments
    \item \textbf{Missing Parameter Errors:} Required parameters omitted
    \item \textbf{Tool Execution Errors:} Runtime failures during tool execution
    \item \textbf{Timeout Errors:} Operations exceeding time limits
    \item \textbf{Result Parsing Errors:} Failures in interpreting tool outputs
    \item \textbf{Result Interpretation Errors:} Incorrect semantic understanding of results
    \item \textbf{Handoff Errors:} Failures in inter-agent communication
    \item \textbf{State Management Errors:} Inconsistencies in shared state variables
    \item \textbf{Retry Logic Errors:} Inappropriate retry behavior
    \item \textbf{Graceful Degradation Errors:} Failures in error recovery paths
\end{enumerate}

\section{Prompt Templates}
\label{app:prompts}

\subsection{RAG System Prompt (Chapter~\ref{ch:rag})}

\begin{verbatim}
You are a helpful AI assistant specialized in answering questions 
about PCI DSS (Payment Card Industry Data Security Standard).

Instructions:
1. Answer questions based ONLY on the provided context
2. If the context doesn't contain relevant information, say so
3. Cite specific sections when possible
4. Be concise but thorough
5. Avoid repeating information unnecessarily

Context:
{retrieved_context}

Question: {user_question}

Answer:
\end{verbatim}

\subsection{Sentiment Analysis System Prompt (Chapter~\ref{ch:strategies})}

\begin{verbatim}
You are an advanced sentiment analysis assistant. Your task is to 
analyze text and provide a sentiment rating along with a brief 
explanation. The sentiment rating should be based on a {scale}-point 
scale: {sentiments}. Always respond with a JSON object containing 
the sentiment and the explanation.

Example format:
- Input: [example text]
- Output: {"sentiment": "[label]", "explanation": "brief reason"}

{few_shot_examples}

- Input: {input_text}
- Output:
\end{verbatim}

\subsection{Logical Reasoning System Prompt (Section~\ref{subsec:turtle_soup})}

\begin{verbatim}
You are an expert in logical reasoning, specifically trained to 
solve Turtle Soup (lateral thinking) puzzles.

Task: Classify the given question-answer pair into one of five 
categories:
- 是 (Yes): The answer confirms the question
- 不是 (No): The answer negates the question
- 不重要 (Unimportant): The detail is not relevant to the puzzle
- 问法错误 (Wrong Question): The question is incorrectly framed
- 回答正确 (Correct Answer): The question reveals the solution

Puzzle Background:
{puzzle_story}

Truth (hidden from players):
{truth}

Question: {question}

Your classification (respond with only the Chinese category):
\end{verbatim}

\subsection{Multi-Agent Data Engineering Prompt (Chapter~\ref{ch:agentic})}

\begin{verbatim}
You are a Data Engineering Agent responsible for database operations
in an invoice reconciliation workflow.

Available Tools:
1. sql_query_builder: Constructs SQL queries from natural language
2. sql_executor: Executes validated SQL queries
3. result_formatter: Formats query results as natural language

Database Schema:
{schema_description}

Instructions:
1. Analyze the user's request
2. Construct appropriate SQL query with bounded results
3. Validate query against schema
4. Execute and format results
5. Return natural language summary
6. Store critical data in state variables (not in message text)

Current State Variables:
{state_variables}

User Request: {request}
\end{verbatim}

\subsection{Soccer Betting Sentiment Prompt (Section~\ref{subsec:soccer_betting})}

\begin{verbatim}
You are an advanced sentiment analysis assistant.
Your task is to analyze match commentaries in the English Premier 
League. The sentiment rating on a certain team should be based on 
a 5-point scale: Strongly Positive (5), Positive (4), Neutral (3), 
Negative (2), or Strongly Negative (1).

I would like you to predict the sentiment for each team for the 
upcoming match based on what happened in the previous match. For 
the upcoming match on [match_date], [home_team] is the home team 
and [away_team] is the away team. I have found the BBC commentaries 
of both teams' previous matches.

Rate the sentiment of each team on an output on a scale of 1-5. 
Reply in a JSON format with the following keys: home_team_sentiment, 
away_team_sentiment. Do not reply with anything else.

Below are the articles:

The article for [home_team]'s previous match on 
[home_team_last_match_date]: [home_team_commentary]

The article for [away_team]'s previous match on 
[away_team_last_match_date]: [away_team_commentary]
\end{verbatim}

\section{Novel Metrics Summary}
\label{app:metrics}

Table~\ref{tab:app_metrics_summary} provides a consolidated reference for all seven novel metrics introduced in this dissertation. These metrics address evaluation gaps identified across the three research pillars, capturing operational dimensions invisible to traditional accuracy-focused assessment.

\begin{table}[htbp]
\centering
\caption{Summary of Novel Evaluation Metrics}
\label{tab:app_metrics_summary}
\resizebox{\textwidth}{!}{%
\begin{tabular}{@{}p{1.2cm}p{4.5cm}p{4cm}p{1.8cm}p{3cm}@{}}
\toprule
\textbf{Metric} & \textbf{Definition} & \textbf{Purpose} & \textbf{Chapter} & \textbf{Key Finding} \\
\midrule
RAP & $P \times (1-\text{RR})^3$ where $P$ = performance, RR = repetition ratio & Penalize repetition in RAG outputs while preserving task performance & Ch.~\ref{ch:rag} & Up to 93.1\% repetition reduction with $<$4\% performance loss \\[1.2em]
VCR & $\dfrac{\text{Valid Classifications}}{\text{Total Outputs}}$ & Measure output reliability for structured classification tasks & Ch.~\ref{ch:strategies} & Fine-tuning improves VCR from 1.17\% to 100\% in 0.2 epochs \\[1.5em]
ECR@1 & $\dfrac{\text{First-Attempt Successes}}{\text{Total Evaluations}}$ & Capture operational reliability and retry cost overhead & Ch.~\ref{ch:agentic} & 85.4\%--100\% range across models; 175$\times$ cost variation \\[1.5em]
ASR & $\dfrac{\sum_{i=1}^{N} s_i}{N}$ where $s_i \in \{0,1\}$ for atomic step $i$ & Fine-grained agentic reliability via N-step decomposition & Ch.~\ref{ch:agentic} & qwen3:32b achieves 98\% ASR matching GPT-4.1 \\[1.5em]
RI & $\dfrac{\text{MRT}}{\text{MRT} + \text{MOT}}$ where MRT = mean reasoning tokens and MOT = mean output tokens & Quantify reasoning-explainability trade-offs & Ch.~\ref{ch:strategies} & Three regimes: Minimal ($<$0.5), Balanced (0.5--0.85), Intensive ($>$0.85) \\[1.5em]
TCR & $\dfrac{\text{Complete Translations}}{\text{Total Entries}}$ & Measure translation reliability (no untranslated segments) & Ch.~\ref{ch:strategies} & Fine-tuning improves TCR to $>$98\% across models \\[1.5em]
CPS & $\dfrac{\text{Source Characters}}{\text{Time (seconds)}}$ & Measure translation efficiency & Ch.~\ref{ch:strategies} & 2--3$\times$ speed improvement with smaller fine-tuned models \\
\bottomrule
\end{tabular}
}
\end{table}

\textbf{Metric Design Principles:}

\begin{enumerate}[leftmargin=*]
    \item \textbf{Complementarity:} Each metric captures a distinct evaluation dimension not addressed by traditional accuracy measures
    \item \textbf{Actionability:} Metrics provide concrete targets for optimization (e.g., RPP tuning for RAP, fine-tuning duration for VCR)
    \item \textbf{Operational Relevance:} Metrics reflect real-world deployment concerns including reliability, efficiency, and cost
    \item \textbf{Cross-Pillar Integration:} Metrics inform decisions across RAG systems (RAP), agentic workflows (ASR, ECR@1), and enterprise deployment (VCR, RI, TCR, CPS)
\end{enumerate}

\textbf{Abbreviations:} RR = Repetition Ratio; MRT = Mean Reasoning Tokens; MOT = Mean Output Tokens; RPP = Repetition Penalty Parameter.